\documentclass[11pt,a4paper]{article}

\usepackage[utf8]{inputenc}
\usepackage[T1]{fontenc}
\usepackage[spanish,es-nodecimaldot]{babel}
\usepackage{lmodern}
\usepackage{microtype}
\usepackage{geometry}
\usepackage{xcolor}
\usepackage{booktabs}
\usepackage{array}
\usepackage{tabularx}
\usepackage{amsmath,amssymb}
\usepackage{siunitx}
\usepackage{hyperref}
\usepackage{cleveref}
\usepackage{enumitem}
\usepackage{graphicx}
\usepackage{tikz}
\usepackage{pgfplots}
\usepackage{listings}
\usepackage{fancyhdr}
\usepackage{titlesec}
\usepackage{caption}

\usetikzlibrary{arrows.meta,positioning,shapes.geometric}
\pgfplotsset{compat=1.18}
\geometry{margin=2.35cm}
\sisetup{output-decimal-marker={.}}

\definecolor{egfnblue}{HTML}{245A73}
\definecolor{egfngreen}{HTML}{2C7A6B}
\definecolor{egfnorange}{HTML}{D97732}
\definecolor{softgray}{HTML}{F3F5F6}
\definecolor{textgray}{HTML}{40464B}

\hypersetup{
  colorlinks=true,
  linkcolor=egfnblue,
  citecolor=egfngreen,
  urlcolor=egfnorange,
  pdftitle={Energy-Gated Frequency Neuron: desarrollo y evaluacion},
  pdfauthor={Henry}
}

\pagestyle{fancy}
\fancyhf{}
\fancyhead[L]{\small Energy-Gated Frequency Neuron}
\fancyhead[R]{\small Reporte técnico}
\fancyfoot[C]{\thepage}
\setlength{\headheight}{14pt}

\titleformat{\section}{\Large\bfseries\color{egfnblue}}{\thesection}{0.65em}{}
\titleformat{\subsection}{\large\bfseries\color{egfngreen}}{\thesubsection}{0.65em}{}
\setlist{itemsep=2pt,topsep=4pt}
\captionsetup{font=small,labelfont=bf}

\lstset{
  basicstyle=\ttfamily\small,
  backgroundcolor=\color{softgray},
  frame=single,
  rulecolor=\color{softgray},
  breaklines=true,
  columns=fullflexible,
  showstringspaces=false
}

\newcommand{\project}{Energy-Gated Frequency Neuron}
\newcommand{\egfn}{EGFN}
\newcommand{\pending}{\textcolor{egfnorange}{\textbf{En curso}}}

\begin{document}

\begin{titlepage}
  \centering
  \vspace*{2.2cm}
  {\color{egfnblue}\rule{\textwidth}{1.2pt}}\\[1.2cm]
  {\Huge\bfseries \project\par}
  \vspace{0.45cm}
  {\LARGE De una neurona frecuencial interpretable a la detección de anomalías industriales\par}
  \vspace{1.2cm}
  {\large Reporte de desarrollo, validación experimental y conclusiones\par}
  \vspace{1.5cm}

  \begin{tikzpicture}[node distance=8mm, every node/.style={font=\small}]
    \node[draw=egfnblue,fill=egfnblue!8,rounded corners,minimum width=2.6cm,minimum height=1cm] (audio) {Audio crudo};
    \node[draw=egfngreen,fill=egfngreen!8,rounded corners,right=of audio,minimum width=2.7cm,minimum height=1cm] (filters) {Filtros};
    \node[draw=egfnorange,fill=egfnorange!10,rounded corners,right=of filters,minimum width=2.7cm,minimum height=1cm] (gates) {Energía + gates};
    \node[draw=egfnblue,fill=egfnblue!8,rounded corners,right=of gates,minimum width=2.8cm,minimum height=1cm] (decision) {Clasificación};
    \draw[-{Latex},thick] (audio) -- (filters);
    \draw[-{Latex},thick] (filters) -- (gates);
    \draw[-{Latex},thick] (gates) -- (decision);
  \end{tikzpicture}

  \vfill
  {\large Autor: Henry\par}
  \vspace{0.25cm}
  {\large Proyecto de investigación y portafolio\par}
  \vspace{0.25cm}
  {\large \today\par}
  \vspace{1cm}
  {\color{egfnblue}\rule{\textwidth}{1.2pt}}
\end{titlepage}

\begin{abstract}
Este proyecto investiga una arquitectura de audio denominada \emph{Energy-Gated
Frequency Neuron} (\egfn). Su hipótesis central es que un banco de filtros,
combinado con medición explícita de energía y compuertas aprendibles, puede
proporcionar un sesgo inductivo útil e interpretable para señales cuyos patrones
discriminativos aparecen en regiones frecuenciales. La investigación evolucionó
desde señales sintéticas controladas hasta voz real y sonidos de maquinaria
industrial. En Google Speech Commands, la variante temporal amplia alcanzó
\(0.886\) de exactitud de prueba frente a \(0.892\) de una Conv1D. En MIMII
Valve obtuvo AUC \(0.984\), comparable y ligeramente superior al \(0.982\) de
la Conv1D. La calibración del umbral elevó el F1 de EGFN de \(0.738\) a
\(0.828\), aunque redujo el recall. Los resultados apoyan el valor del frontend
frecuencial, pero también muestran que los gates actuales no explican por sí
solos las decisiones y que los filtros no restringidos pueden perder su forma
clásica de banda. Se presenta una evaluación honesta de estas fortalezas,
limitaciones y próximos experimentos.
\end{abstract}

\tableofcontents
\newpage

\section{Origen y motivación}

El proyecto surgió de una pregunta sencilla: ¿puede una neurona diseñada con
principios de procesamiento de señales identificar qué regiones de frecuencia
son relevantes antes de tomar una decisión? Una red Conv1D convencional puede
aprender filtros directamente desde el audio, pero sus canales no tienen por
qué conservar una interpretación física inmediata. La propuesta EGFN introduce
explícitamente tres elementos:

\begin{enumerate}
  \item un banco de filtros inicializado con bandas conocidas;
  \item una medida de energía por banda;
  \item una compuerta aprendible que modula la respuesta de cada banda.
\end{enumerate}

El objetivo no fue afirmar desde el inicio que EGFN superaría universalmente a
una CNN, sino comprobar si podía aproximarse a su capacidad predictiva mientras
ofrecía una representación frecuencial más fácil de examinar. Esta diferencia
es especialmente relevante en vibración, maquinaria, voz y audio médico, donde
los cambios de estado suelen expresarse como variaciones espectrales.

\section{Hipótesis y preguntas de investigación}

La hipótesis principal es:

\begin{quote}
Un frontend neuronal con estructura frecuencial explícita puede alcanzar una
capacidad discriminativa comparable a un baseline convolucional y, al mismo
tiempo, exponer información útil sobre las regiones espectrales procesadas.
\end{quote}

De ella se derivan cuatro preguntas:

\begin{enumerate}
  \item ¿Aprende el bloque EGFN en señales donde la clase depende de frecuencia?
  \item ¿Generaliza desde datos sintéticos hasta audio real con variación de hablante?
  \item ¿Puede competir con Conv1D en una tarea industrial normal/anormal?
  \item ¿Los filtros, energías y gates proporcionan evidencia interpretable?
\end{enumerate}

\section{Formulación de la neurona frecuencial}

Sea (x[n]) una señal de audio y (h_k[n]) el filtro asociado con la banda
\(k\). La respuesta filtrada es

\begin{equation}
u_k[n] = (h_k * x)[n].
\end{equation}

Después se aplica un sesgo y una no linealidad \(\phi\):

\begin{equation}
a_k[n] = \phi\!\left(u_k[n] + b_k\right).
\end{equation}

La energía media de la banda se calcula como

\begin{equation}
E_k = \frac{1}{N}\sum_{n=1}^{N} u_k[n]^2.
\end{equation}

En el modo independiente, el gate es

\begin{equation}
g_k = \sigma\!\left(\alpha_k \log(1+E_k)+\beta_k\right),
\end{equation}

y la salida modulada queda

\begin{equation}
y_k[n] = g_k a_k[n].
\end{equation}

También se implementó un modo contextual en el cual el vector completo de
energías entra a una MLP:

\begin{equation}
\mathbf{g}=\sigma\!\left(\operatorname{MLP}
(\log(1+\mathbf{E}))\right).
\end{equation}

Este modo permite que la activación de una banda dependa del patrón energético
global. En los experimentos finales se utilizó el modo independiente porque
ofreció una comparación más directa y conservó parámetros por banda.

\section{Evolución arquitectónica}

\subsection{Clasificador EGFN inicial}

La primera versión resumía cada banda mediante media, máximo, desviación
estándar, energía y gate. Estas características entraban a un pequeño MLP. Era
una arquitectura compacta y adecuada para validar el bloque, pero el pooling
global eliminaba gran parte del orden temporal.

\subsection{EGFN temporal}

Para preservar patrones locales se agregó una cabeza convolucional temporal:

\begin{center}
\begin{tikzpicture}[
  node distance=6mm,
  block/.style={draw,rounded corners,minimum height=8mm,minimum width=2.2cm,align=center,font=\small},
  arrow/.style={-{Latex},thick}
]
  \node[block,fill=egfnblue!8] (x) {Waveform\\\(1\times N\)};
  \node[block,fill=egfngreen!10,right=of x] (bank) {Banco EGFN\\8 bandas};
  \node[block,fill=egfnorange!10,right=of bank] (gate) {Energía y\\gating};
  \node[block,fill=egfnblue!8,right=of gate] (proj) {Proyección\\\(1\times1\)};
  \node[block,fill=egfngreen!10,below=of proj] (temp) {Conv temporal\\64, 128};
  \node[block,fill=egfnorange!10,left=of temp] (pool) {Pooling\\global};
  \node[block,fill=egfnblue!8,left=of pool] (out) {Logits\\de clase};
  \draw[arrow] (x) -- (bank);
  \draw[arrow] (bank) -- (gate);
  \draw[arrow] (gate) -- (proj);
  \draw[arrow] (proj) -- (temp);
  \draw[arrow] (temp) -- (pool);
  \draw[arrow] (pool) -- (out);
\end{tikzpicture}
\end{center}

La variante \emph{wide} proyecta las ocho bandas a 32 canales y usa bloques
temporales de 64 y 128 canales. Incluye Batch Normalization, GELU, pooling,
dropout y clasificación final. Los filtros pueden mantenerse fijos o hacerse
aprendibles; la mejor variante utilizó filtros aprendibles.

\subsection{Baseline Conv1D}

Conv1D sirve como control experimental. Aprende directamente kernels sobre la
forma de onda y representa una alternativa estándar con menor compromiso de
interpretabilidad explícita. Una comparación válida exige mismo dataset,
split, augmentación y evaluación.

\section{Ruta experimental}

\subsection{Señales sintéticas}

La primera etapa generó tonos y combinaciones de frecuencia con ruido
controlado. Su función fue verificar formas tensoriales, propagación hacia
adelante y capacidad de aprendizaje. En la tarea sencilla, el modelo llegó a
exactitud perfecta. Al aumentar solapamiento, ruido y dificultad, la exactitud
bajó al rango aproximado de \(0.70\)--\(0.75\), mostrando que el primer éxito no
debía interpretarse como evidencia suficiente de generalización.

El barrido de SNR de una variante difícil produjo resultados alrededor de
\(0.69\)--\(0.75\), incluso con ruido intenso. Esta fase permitió comparar
filtros fijos, filtros aprendibles, gating contextual y regularización de gates.

\subsection{Free Spoken Digit Dataset}

FSDD fue el primer contacto con audio real. Un split aleatorio alcanzó
\(0.840\) de exactitud de prueba, mientras que splits por hablante fueron más
difíciles. La variante temporal alcanzó \(0.724\) de mejor validación y
\(0.708\) en test en uno de estos escenarios. La brecha entre entrenamiento y
validación reveló sobreajuste y dependencia de identidad vocal. Este resultado
motivó el salto a un dataset más grande y diverso.

\subsection{Google Speech Commands}

Se utilizaron diez comandos: \texttt{yes}, \texttt{no}, \texttt{up},
\texttt{down}, \texttt{left}, \texttt{right}, \texttt{on}, \texttt{off},
\texttt{stop} y \texttt{go}. El split oficial contenía 30,769 ejemplos de
entrenamiento, 3,703 de validación y 4,074 de prueba.

\begin{table}[ht]
\centering
\caption{Resultados en Google Speech Commands.}
\label{tab:speech}
\begin{tabular}{lccc}
\toprule
Modelo & Mejor validación & Test accuracy & Gate medio \\
\midrule
Conv1D & \textbf{0.902} & \textbf{0.892} & -- \\
EGFN temporal & 0.856 & 0.849 & 0.547 \\
EGFN temporal wide & 0.893 & 0.886 & 0.444 \\
\bottomrule
\end{tabular}
\end{table}

La diferencia final entre Conv1D y EGFN wide fue de solo \(0.006\). Esto
demostró que el frontend frecuencial podía aproximarse al baseline en una tarea
real y amplia, justificando una evaluación donde la estructura espectral fuera
central.

\section{Aplicación industrial con MIMII}

\subsection{Dataset y tarea}

MIMII contiene sonidos normales y anormales de maquinaria industrial
\cite{purohit2019mimii}. Se descargó el subconjunto de válvulas a \(6\,\mathrm{dB}\),
con 4,170 grabaciones: 2,922 para entrenamiento, 624 para validación y 624
para test. La tarea actual es clasificación binaria supervisada.

Existe un fuerte desbalance: las grabaciones normales son mucho más numerosas.
Por ello se usó cross-entropy ponderada y se reportaron precision, recall, F1 y
AUC, además de accuracy. Se aplicaron augmentaciones de ganancia, desplazamiento
temporal, ruido y enmascaramiento.

\subsection{Resultados con umbral 0.50}

\begin{table}[ht]
\centering
\caption{Resultados de prueba en MIMII Valve con umbral 0.50.}
\label{tab:mimii-default}
\begin{tabular}{lccccc}
\toprule
Modelo & Accuracy & Precision & Recall & F1 & AUC \\
\midrule
Conv1D & \textbf{0.955} & \textbf{0.765} & 0.873 & \textbf{0.816} & 0.982 \\
EGFN temporal wide & 0.929 & 0.639 & 0.873 & 0.738 & \textbf{0.984} \\
\bottomrule
\end{tabular}
\end{table}

Conv1D produjo mejor accuracy, precision y F1 en el punto operativo estándar.
Sin embargo, EGFN igualó exactamente el recall y obtuvo un AUC ligeramente
superior. El AUC indica que la capacidad de ordenar ejemplos normales y
anormales es prácticamente equivalente; la diferencia principal estaba en el
umbral usado para convertir probabilidad en decisión.

\section{Calibración del umbral}

El umbral se seleccionó maximizando F1 únicamente sobre validación. Después se
aplicó una sola vez al test, evitando optimizar con información del conjunto
final. Los umbrales seleccionados fueron \(0.65\) para Conv1D y \(0.78\) para
EGFN.

\begin{table}[ht]
\centering
\caption{Comparación de puntos operativos calibrados.}
\label{tab:calibrated}
\resizebox{\textwidth}{!}{%
\begin{tabular}{llcccccc}
\toprule
Modelo & Umbral & Accuracy & Precision & Recall & F1 & FP & FN \\
\midrule
Conv1D & 0.50 & 0.955 & 0.765 & 0.873 & 0.816 & 19 & 9 \\
Conv1D & 0.65 & 0.960 & 0.883 & 0.746 & 0.809 & 7 & 18 \\
EGFN & 0.50 & 0.929 & 0.639 & 0.873 & 0.738 & 35 & 9 \\
EGFN & 0.78 & \textbf{0.965} & \textbf{0.930} & 0.746 & \textbf{0.828} & \textbf{4} & 18 \\
\bottomrule
\end{tabular}}
\end{table}

La calibración de EGFN redujo falsas alarmas de 35 a 4 y elevó F1 de 0.738 a
0.828. El costo fue aumentar anomalías omitidas de 9 a 18. Por tanto, no existe
un umbral universalmente mejor:

\begin{itemize}
  \item \textbf{Umbral 0.50:} preferible cuando omitir una falla es muy costoso.
  \item \textbf{Umbral 0.78:} preferible cuando las falsas alarmas son costosas.
\end{itemize}

\section{Análisis de interpretabilidad}

\subsection{Gates y energía}

Los gates promedio se mantuvieron alrededor de \(0.50\)--\(0.56\) y mostraron
diferencias pequeñas entre audio normal y anormal. Por ejemplo, el mayor cambio
medio observado fue aproximadamente \(-0.012\) en la banda inicial de
1.5--2.5 kHz. Esto significa que no es correcto afirmar que el gate por sí solo
explica la clasificación.

Las energías sí mostraron diferencias más visibles. El audio anormal presentó
menor energía media en la banda inicial de 1.5--2.5 kHz y mayor energía sobre
4 kHz. La cabeza temporal recibe las respuestas moduladas completas, por lo que
puede explotar tanto energía como estructura temporal aunque los gates medios
sean similares.

\subsection{Filtros aprendibles}

Los kernels se inicializaron como filtros pasa-baja, pasa-banda y pasa-alta,
pero se entrenaron sin una restricción que conservara esas formas. El análisis
FFT mostró desplazamientos considerables y respuestas multibanda. Por ejemplo,
el filtro inicialmente menor de 250 Hz terminó con un pico aproximado en
1.90 kHz, mientras que el filtro inicialmente mayor de 4 kHz alcanzó un pico
cercano a 6.34 kHz.

Este hallazgo tiene dos interpretaciones:

\begin{enumerate}
  \item \textbf{Ventaja:} el modelo puede adaptar libremente su representación
  al sonido de las válvulas.
  \item \textbf{Limitación:} la etiqueta de banda inicial deja de describir
  exactamente la respuesta final y disminuye la interpretabilidad clásica.
\end{enumerate}

Una futura versión puede parametrizar directamente frecuencias de corte o usar
filtros tipo SincNet \cite{ravanelli2018sincnet}, conservando bandas válidas
durante el aprendizaje.

\section{Comparación conceptual con otras arquitecturas}

\begin{table}[ht]
\centering
\caption{Comparación cualitativa de enfoques de audio.}
\label{tab:architectures}
\begin{tabularx}{\textwidth}{>{\bfseries}lXXX}
\toprule
Arquitectura & Representación & Ventaja principal & Limitación principal \\
\midrule
Conv1D & Filtros libres sobre waveform & Buen rendimiento y simplicidad & Canales menos interpretables \\
CNN sobre espectrograma & Imagen tiempo-frecuencia & Patrones espectro-temporales explícitos & Requiere transformación previa \\
MFCC + MLP & Coeficientes acústicos compactos & Bajo costo computacional & Descarta detalles del waveform \\
SincNet & Cortes frecuenciales parametrizados & Filtros físicamente restringidos & Menor libertad que kernels completos \\
EGFN temporal & Filtros, energía, gates y Conv temporal & Sesgo frecuencial y diagnóstico por banda & Gates poco selectivos y filtros no restringidos \\
\bottomrule
\end{tabularx}
\end{table}

EGFN no debe presentarse como reemplazo universal de Conv1D. Su contribución es
un frontend híbrido donde parte del procesamiento tiene significado espectral
explícito. El resultado más convincente es alcanzar discriminación comparable
en MIMII, no una superioridad absoluta en todas las métricas.

\section{Experimento multisemilla}

\pending. Una sola semilla puede favorecer accidentalmente a un modelo por la
inicialización, el split o el orden de augmentación. Para medir estabilidad se
diseñó un runner reproducible con semillas \(42\), \(123\) y \(456\) para
Conv1D y EGFN. Cada corrida:

\begin{enumerate}
  \item usa la misma semilla para Python, NumPy, PyTorch y CUDA;
  \item construye splits emparejados por semilla;
  \item entrena con la configuración establecida;
  \item selecciona su threshold usando únicamente validación;
  \item reporta métricas default y calibradas sobre test.
\end{enumerate}

El resultado final se expresará como

\begin{equation}
\bar{x} \pm s,
\qquad
s=\sqrt{\frac{1}{n-1}\sum_{i=1}^{n}(x_i-\bar{x})^2},
\end{equation}

donde \(s\) es la desviación estándar muestral y \(n=3\). La
\cref{tab:multiseed} queda preparada para completar al finalizar el proceso.

\begin{table}[ht]
\centering
\caption{Resultados multisemilla pendientes, media \(\pm\) desviación estándar.}
\label{tab:multiseed}
\begin{tabular}{llccc}
\toprule
Modelo & Punto operativo & Accuracy & F1 & AUC \\
\midrule
Conv1D & 0.50 & \multicolumn{3}{c}{En ejecución} \\
Conv1D & Calibrado & \multicolumn{3}{c}{En ejecución} \\
EGFN & 0.50 & \multicolumn{3}{c}{En ejecución} \\
EGFN & Calibrado & \multicolumn{3}{c}{En ejecución} \\
\bottomrule
\end{tabular}
\end{table}

\section{Validez, limitaciones y riesgos}

\begin{enumerate}
  \item \textbf{Split de MIMII.} El experimento actual usa un split estratificado
  de grabaciones y mezcla los mismos machine IDs entre conjuntos. Esto mide
  generalización a nuevas grabaciones, no necesariamente a una máquina física
  nunca observada.
  \item \textbf{Clasificación supervisada.} Se usan ejemplos anormales durante
  entrenamiento. Un sistema clásico de detección de anomalías puede entrenarse
  solo con normalidad; es una tarea distinta y más cercana a escenarios donde
  las fallas son escasas.
  \item \textbf{Desbalance.} Accuracy puede ocultar errores en la clase
  minoritaria. Por eso AUC, recall, precision y F1 son indispensables.
  \item \textbf{Interpretabilidad parcial.} Los gates son observables, pero su
  existencia no demuestra causalidad. Se requieren ablaciones o intervenciones
  por banda para sostener explicaciones más fuertes.
  \item \textbf{Filtros no restringidos.} Los kernels aprendidos pueden dejar de
  comportarse como filtros clásicos.
  \item \textbf{Escala experimental.} Tres semillas permiten una primera medida
  de variabilidad, pero no equivalen a validación industrial extensa.
\end{enumerate}

\section{Conclusiones}

El proyecto cumplió su objetivo principal como investigación aplicada y pieza
de portafolio. Se diseñó un módulo PyTorch propio, se verificó en señales
sintéticas, se evaluó en dos datasets de voz y se transfirió a detección de
fallas industriales. La variante EGFN temporal wide quedó a 0.006 de Conv1D en
Speech Commands y obtuvo AUC 0.984 en MIMII Valve frente a 0.982 del baseline.

El análisis posterior fue tan importante como la métrica principal. La
calibración mostró que EGFN puede operar con muy pocas falsas alarmas, pero a
costa de recall. El análisis frecuencial mostró cambios de energía útiles y, al
mismo tiempo, reveló que los gates son poco discriminativos y que los filtros
libres pueden perder su interpretación inicial.

La conclusión defendible no es que EGFN sea universalmente superior, sino que
\textbf{alcanza capacidad discriminativa comparable a Conv1D en una aplicación
industrial mientras expone una representación frecuencial analizable}. Esta es
una base sólida para continuar con filtros restringidos, evaluación por
machine ID, detección normal-only y despliegue de inferencia.

\section{Próximos pasos}

\begin{enumerate}
  \item Completar y reportar el experimento multisemilla.
  \item Evaluar generalización dejando un machine ID fuera del entrenamiento.
  \item Implementar filtros parametrizados por cortes frecuenciales válidos.
  \item Realizar ablaciones: sin gates, filtros fijos y energía sin gating.
  \item Medir parámetros, latencia, memoria e inferencia en GPU y CPU.
  \item Preparar README final, licencia, entorno reproducible y release v1.0.
  \item Como extensión, probar bombas, ventiladores o audio respiratorio.
\end{enumerate}

\appendix
\section{Comandos principales}

\begin{lstlisting}[language=bash,caption={Entrenamiento del baseline industrial.}]
python scripts/train_mimii.py --data-dir data/raw/mimii \
  --machine-type valve --model conv1d --augment \
  --balanced-loss --num-workers 2 --epochs 40 \
  --output-dir outputs/mimii_valve_conv1d
\end{lstlisting}

\begin{lstlisting}[language=bash,caption={Entrenamiento de EGFN temporal wide.}]
python scripts/train_mimii.py --data-dir data/raw/mimii \
  --machine-type valve --model egfn_temporal \
  --filter-bank fine --learnable-filters \
  --gate-mode independent --frontend-channels 32 \
  --temporal-channels 64,128 --dropout 0.25 \
  --label-smoothing 0.05 --scheduler cosine \
  --augment --balanced-loss --num-workers 2 \
  --epochs 40 --patience 10 \
  --output-dir outputs/mimii_valve_egfn_temporal_wide
\end{lstlisting}

\begin{lstlisting}[language=bash,caption={Calibración e interpretabilidad.}]
python scripts/analyze_mimii_results.py \
  --run-dir outputs/mimii_valve_egfn_temporal_wide \
  --num-workers 2 --device cuda
\end{lstlisting}

\begin{lstlisting}[language=bash,caption={Experimento multisemilla reproducible.}]
python scripts/run_mimii_multiseed.py \
  --device cuda --num-workers 2
\end{lstlisting}

\begin{thebibliography}{9}

\bibitem{purohit2019mimii}
H.~Purohit, R.~Tanabe, T.~Ichige, T.~Endo, Y.~Nikaido, K.~Suefusa y Y.~Kawaguchi.
\newblock MIMII Dataset: Sound Dataset for Malfunctioning Industrial Machine Investigation and Inspection.
\newblock \emph{Proceedings of the Detection and Classification of Acoustic Scenes and Events Workshop}, 2019.

\bibitem{warden2018speech}
P.~Warden.
\newblock Speech Commands: A Dataset for Limited-Vocabulary Speech Recognition.
\newblock \emph{arXiv preprint arXiv:1804.03209}, 2018.

\bibitem{ravanelli2018sincnet}
M.~Ravanelli y Y.~Bengio.
\newblock Speaker Recognition from Raw Waveform with SincNet.
\newblock \emph{IEEE Spoken Language Technology Workshop}, 2018.

\end{thebibliography}

\end{document}
